\documentclass[11pt]{article}
\usepackage{amsmath,amssymb,amsthm}
\usepackage{hyperref}

\newtheorem{theorem}{Theorem}
\newtheorem{proposition}[theorem]{Proposition}
\newtheorem{lemma}[theorem]{Lemma}
\newtheorem{conjecture}[theorem]{Conjecture}
\newtheorem{definition}[theorem]{Definition}
\newtheorem{remark}[theorem]{Remark}
\newtheorem{example}[theorem]{Example}
\newtheorem{observation}[theorem]{Observation}

\newcommand{\N}{\mathbb{N}}
\newcommand{\Z}{\mathbb{Z}}
\newcommand{\Q}{\mathbb{Q}}
\newcommand{\CPP}{\operatorname{CPP}}
\newcommand{\RPP}{\operatorname{RPP}}
\newcommand{\SSYT}{\operatorname{SSYT}}
\newcommand{\GK}{\operatorname{GK}}
\newcommand{\QSym}{\operatorname{QSym}}
\newcommand{\Sym}{\operatorname{Sym}}

% q-series notation
\newcommand{\qpoch}[2]{(#1;#2)_\infty}
\newcommand{\qpochn}[3]{(#1;#2)_{#3}}
\newcommand{\qbin}[2]{\genfrac{[}{]}{0pt}{}{#1}{#2}}

\title{The maximum entry in reverse plane partitions\\
and cylindric plane partitions}
\author{Robin Langer}
\date{\today}

\begin{document}
\maketitle

\begin{abstract}
We collect what is known about reverse plane partitions (RPPs)
and cylindric plane partitions (CPPs) refined by their maximum
entry. The maximum entry~$n$ governs the ``depth'' of the
combinatorial object: it is the number of layers in the
level-set decomposition, the number of non-intersecting lattice
paths in the NILP model, and the number of insertion steps in
the RSK/growth-diagram framework. We describe the transfer
matrix formulation (products of skew Schur functions),
Borodin's product formula, the E--H factorization that extracts
Warnaar's $Q_{n,c}(q)$ polynomials, and the partial
quasi-symmetry of the frequency refinement at fixed
maximum entry.
\end{abstract}

\tableofcontents

%% ========================================================================
\section{Reverse plane partitions and the maximum entry}
\label{sec:rpp}
%% ========================================================================

\subsection{Definitions}

A \emph{reverse plane partition} (RPP) of shape~$\lambda$ is a
filling of the Young diagram of~$\lambda$ with nonneg integers
that is weakly increasing along rows (left to right) and weakly
increasing down columns (top to bottom). We write $\RPP(\lambda)$
for the set of all RPPs of shape~$\lambda$.

\begin{remark}
Some authors use the convention that entries are weakly
\emph{decreasing} (``plane partitions'' in the sense of
MacMahon). We follow the reverse/increasing convention, which
matches the interlacing-sequence model and Stanley's
$P$-partition framework.
\end{remark}

The \emph{weight} of an RPP~$\pi$ is
$|\pi| = \sum_{(i,j) \in \lambda} \pi_{ij}$,
and the \emph{maximum entry} is
$\max(\pi) = \max_{(i,j)} \pi_{ij}$.

Stanley's celebrated formula~\cite{stanley-1972} gives the
weight generating function:
\begin{equation}\label{eq:stanley-hook}
  \sum_{\pi \in \RPP(\lambda)} q^{|\pi|}
  = \prod_{u \in \lambda} \frac{1}{1 - q^{h_\lambda(u)}},
\end{equation}
where $h_\lambda(u)$ is the hook length at box~$u$.

\subsection{Level-set decomposition}

Given an RPP~$\pi$ with $\max(\pi) = n$, define the
\emph{level sets}:
\[
  L_k(\pi) = \{(i,j) \in \lambda : \pi_{ij} \geq k\},
  \qquad k = 1, 2, \ldots, n.
\]
Since $\pi$ is weakly increasing in both directions,
each $L_k(\pi)$ is a Young diagram (the upper-left corner
condition is preserved). We obtain a chain of partitions:
\begin{equation}\label{eq:level-chain}
  \lambda = L_0 \supseteq L_1 \supseteq L_2
  \supseteq \cdots \supseteq L_n \supseteq \emptyset,
\end{equation}
and the weight decomposes as
$|\pi| = \sum_{k=1}^{n} |L_k|$
(each box at height~$\geq k$ contributes 1 for each level).

The RPP is uniquely determined by its level-set chain.
Thus:

\begin{proposition}\label{prop:rpp-level-sets}
RPPs of shape $\lambda$ with $\max \leq n$ are in bijection
with chains of partitions
$\lambda \supseteq \mu^{(1)} \supseteq \cdots \supseteq
\mu^{(n)} \supseteq \emptyset$
inside~$\lambda$.
\end{proposition}

This is the ``$n$-layer'' decomposition: max entry~$= n$
means exactly $n$ layers.

\subsection{The interlacing-sequence model}

Reading an RPP along its anti-diagonals gives an
\emph{interlacing sequence} of partitions. For shape~$\lambda$
with profile (binary string) $\pi = (\pi_1, \ldots, \pi_T)$:
\[
  \emptyset = \nu^{(0)},\; \nu^{(1)},\; \ldots,\;
  \nu^{(T)},\; \nu^{(T+1)} = \emptyset,
\]
where $\nu^{(k-1)}$ and $\nu^{(k)}$ satisfy the interlacing
condition determined by~$\pi_k$:
\begin{itemize}
  \item If $\pi_k = 0$: $\nu^{(k)}/\nu^{(k-1)}$ is a
    horizontal strip (row added);
  \item If $\pi_k = 1$: $\nu^{(k-1)}/\nu^{(k)}$ is a
    horizontal strip (row removed).
\end{itemize}

In this model, the maximum entry of the RPP equals
$\max_{k,j} \nu^{(k)}_j$, the maximum part across all
partitions in the sequence.

\subsection{Bounding the maximum entry}
\label{sec:rpp-bounded}

For RPPs of shape $\lambda$ with max entry $\leq m$, we
restrict all parts $\nu^{(k)}_j \leq m$ in the interlacing
sequence. This is equivalent to requiring that each partition
$\nu^{(k)}$ has at most $m$ columns (in English convention),
or equivalently that the conjugate $(\nu^{(k)})'$ has at most
$m$ parts.

The generating function for RPPs of shape $\lambda$ refined
by maximum entry is:
\begin{equation}\label{eq:rpp-max-gf}
  \sum_{\pi \in \RPP(\lambda)} z^{\max(\pi)} q^{|\pi|}
  = \sum_{n \geq 0} g_{n,\lambda}(q)\, z^n.
\end{equation}
The variable~$z$ tracks the maximum entry, and
$g_{n,\lambda}(q)$ counts RPPs of shape~$\lambda$ with exact
max entry~$n$ by weight.

%% ========================================================================
\section{Cylindric plane partitions}
\label{sec:cpp}
%% ========================================================================

\subsection{Definition and generating function}

A \emph{rank-$k$ cylindric plane partition} with profile
$c = (c_1, \ldots, c_k)$ is a $k$-tuple of partitions
$(\nu^{(1)}, \ldots, \nu^{(k)})$ satisfying the
\emph{cyclic} interlacing conditions:
\begin{equation}\label{eq:cpp-interlacing}
  \nu^{(i)}_j \geq \nu^{(i+1)}_{j + c_{i+1}},
  \qquad i = 1, \ldots, k,
  \quad \text{indices mod $k$}.
\end{equation}
When $k = 3$ and $c = (c_0, c_1, c_2)$ with $d = |c|$:
\[
  \nu^{(1)}_j \geq \nu^{(2)}_{j+c_1}, \qquad
  \nu^{(2)}_j \geq \nu^{(3)}_{j+c_2}, \qquad
  \nu^{(3)}_j \geq \nu^{(1)}_{j+c_0}.
\]

RPPs are the special case where the profile corresponds
to a partition shape and the boundary condition is
$\nu^{(0)} = \nu^{(T+1)} = \emptyset$ (non-cyclic).
The cyclic boundary condition is what distinguishes CPPs.

The two-variable generating function is:
\begin{equation}\label{eq:GK}
  \GK_c(z, q) = \sum_{\pi \in \CPP(c)} z^{n(\pi)} q^{|\pi|}
  = \sum_{n \geq 0} g_n(q)\, z^n,
\end{equation}
where $n(\pi) = \max_{i,j} \nu^{(i)}_j$ is the maximum entry.

\subsection{Borodin's product formula}

Borodin~\cite{borodin-2007} proved:
\begin{equation}\label{eq:borodin}
  \GK_c(1, q) = \sum_{\pi \in \CPP(c)} q^{|\pi|}
  = \prod_{s \in \mathcal{C}(c)} \frac{1}{1 - q^{h(s)}},
\end{equation}
where $\mathcal{C}(c)$ is the ``cylindric shape'' determined
by the profile~$c$, and $h(s)$ is the hook length of box~$s$
in this cylindric shape. This generalizes
Stanley's formula~\eqref{eq:stanley-hook} for RPPs.

The $z$-refined version involves the transfer matrix
(Section~\ref{sec:transfer-matrix}).

%% ========================================================================
\section{The NILP model: max entry = number of paths}
\label{sec:nilp}
%% ========================================================================

In the non-intersecting lattice path (NILP) interpretation:
\begin{quote}
A CPP with max entry $= n$ corresponds to $n$ non-intersecting
lattice paths on the cylinder.
\end{quote}

Each ``layer'' (the level set $\{(i,j) : \nu^{(i)}_j \geq k\}$)
contributes one path. The paths are non-intersecting because
the interlacing conditions force the level sets to be nested.

The maximum entry thus controls the \emph{combinatorial
complexity}: $n = 1$ gives a single path (equivalent to a
single interlacing chain), while larger~$n$ requires
coordinating $n$ nested paths on the cylinder.

For the positivity conjecture, the case $n = 1$ is trivial
($g_1$ counts objects directly), and the difficulty grows
with~$n$ because the alternating-sign E--H extraction
(Section~\ref{sec:EH}) involves deeper cancellation.

%% ========================================================================
\section{The transfer matrix and Schur functions}
\label{sec:transfer-matrix}
%% ========================================================================

\subsection{Vertex operator formulation}

The generating function $\GK_c(z,q)$ can be expressed using
vertex operators acting on the ring of symmetric functions.
The Pieri rules for Schur functions give:
\[
  s_\mu[X] \cdot h_r[X]
  = \sum_{\lambda \in U_r(\mu)} s_\lambda[X],
  \qquad
  s_\lambda[X + w]
  = \sum_{\mu \in D_r(\lambda)} s_\mu[X]\, w^r,
\]
where $U_r(\mu)$ (resp.\ $D_r(\lambda)$) is the set of
partitions obtained from $\mu$ by adding (resp.\ removing) a
horizontal $r$-strip.

The coefficient of $z^n$ in $\GK_c(z,q)$ is:
\begin{equation}\label{eq:transfer-matrix}
  g_n(q) = \sum_{\vec{\lambda}}
  \prod_{i=1}^{k}
  s_{\lambda^{(i+1)}/\lambda^{(i)}}(1, q, q^2, \ldots),
\end{equation}
where the sum is over $k$-tuples of partitions
$\vec{\lambda} = (\lambda^{(1)}, \ldots, \lambda^{(k)})$
satisfying the cyclic interlacing~\eqref{eq:cpp-interlacing}
with all parts~$\leq n$.

\subsection{Schur positivity}

Each factor $s_{\mu/\nu}(1, q, q^2, \ldots)$ in
\eqref{eq:transfer-matrix} is the principal specialization
of a skew Schur function, which counts semistandard Young
tableaux of skew shape~$\mu/\nu$ by major index:
\[
  s_{\mu/\nu}(1, q, q^2, \ldots)
  = \sum_{T \in \SSYT(\mu/\nu)} q^{|T|}.
\]
Before specialization, \eqref{eq:transfer-matrix} is a sum
of products of skew Schur functions. Products of skew Schur
functions are Schur-positive by the Littlewood--Richardson
rule, so:

\begin{proposition}\label{prop:schur-pos}
The ``unspecialized'' transfer matrix expression
\[
  \widetilde{g}_n^{\mathrm{TM}}(X)
  = \sum_{\vec{\lambda}}
  \prod_{i=1}^{k}
  s_{\lambda^{(i+1)}/\lambda^{(i)}}(X)
\]
is a Schur-positive symmetric function.
Under the principal specialization $x_j \mapsto q^{j-1}$,
it recovers $g_n(q)$.
\end{proposition}

\begin{remark}
This ``unspecialized'' lift tracks a \emph{single} alphabet
$X$ shared across all $k$ skew Schur factors. It refines
the total weight but does \emph{not} distinguish which
factor (= which diagonal of the CPP) a box belongs to.
The full frequency refinement
$\tilde{g}_n(z_1, z_2, \ldots)$ from
Section~\ref{sec:frequency} is a different (and finer)
multivariate lift, and it is \emph{not} symmetric.
\end{remark}

\subsection{Connection to SSYT}

For RPPs of shape $\lambda$ (the non-cylindric case),
the transfer matrix reduces to:
\[
  g_{n,\lambda}^{\mathrm{RPP}}(q)
  = \sum_{\substack{\mu^{(0)} = \emptyset \subseteq
    \mu^{(1)} \subseteq \cdots \subseteq \mu^{(n)} \\
    \mu^{(n)} = \lambda}}
  \prod_{i=1}^{n} s_{\mu^{(i)}/\mu^{(i-1)}}(1, q, q^2, \ldots).
\]
The unspecialized version $\sum \prod_i s_{\mu^{(i)}/\mu^{(i-1)}}(X)$
is a single Schur function evaluated at a product of
alphabets: it equals $s_\lambda(X^n)$ under appropriate
plethystic substitution.

More concretely, RPPs of shape~$\lambda$ with
$\max \leq n$ are in bijection with semistandard Young
tableaux of shape~$\lambda$ with entries in $\{1, \ldots, n\}$
(via the level-set/Gelfand--Tsetlin pattern correspondence),
giving:
\begin{equation}\label{eq:rpp-schur}
  \sum_{\substack{\pi \in \RPP(\lambda) \\
  \max(\pi) \leq n}} q^{|\pi|}
  = s_\lambda(1, q, q^2, \ldots, q^{n-1}).
\end{equation}

The \emph{exact} max-entry count $g_{n,\lambda}(q)$ from
\eqref{eq:rpp-max-gf} is then recovered by differencing:
\begin{equation}\label{eq:rpp-exact-max}
  g_{n,\lambda}(q)
  = s_\lambda(1, q, \ldots, q^{n-1})
  - s_\lambda(1, q, \ldots, q^{n-2}),
\end{equation}
the contribution of RPPs whose maximum entry is exactly~$n$.
This is the RPP analogue of the E--H extraction: the
cumulative $\leq n$ formula is clean (a single Schur
specialization), but isolating the $= n$ piece requires
a subtraction that destroys manifest positivity.

\begin{remark}
For RPPs, positivity of~\eqref{eq:rpp-exact-max} is
immediate: $g_{n,\lambda}(q) \geq 0$ because it counts
objects. But the analogous extraction for CPPs ---
Warnaar's $Q_{n,c}(q)$ in~\eqref{eq:Q-def} --- involves
the infinite product $(zq;q)_\infty$ rather than a simple
difference, making positivity a nontrivial conjecture.
\end{remark}

%% ========================================================================
\section{Growth diagrams and RSK}
\label{sec:rsk}
%% ========================================================================

\subsection{Fomin's growth diagrams}

Fomin's theory of growth diagrams~\cite{fomin-1986,fomin-1995}
provides a unified framework for RSK-type correspondences.
A \emph{growth diagram} is a two-dimensional grid of
partitions satisfying local rules at each cell.

For an RPP $\pi$ of shape $\lambda$ with max entry~$n$:
\begin{enumerate}
  \item The RPP is encoded as an $n$-step process,
    where step~$k$ records the $k$-th level set~$L_k$.
  \item Each step corresponds to one ``row'' of the growth
    diagram.
  \item The local rules (Burge, Knuth, or dual) govern
    how partitions grow at each cell.
  \item The boundary conditions encode the shape~$\lambda$
    and the level-set structure.
\end{enumerate}

The key observation:
\begin{quote}
\emph{The maximum entry of the RPP equals the number of rows
(= insertion steps) in the growth diagram.}
\end{quote}

\subsection{RSK and the maximum entry}

Under the RSK correspondence (in the growth-diagram
formulation):

\begin{itemize}
  \item An RPP~$\pi$ of shape~$\lambda$ with $\max(\pi) = n$
    maps to a pair $(P, Q)$ of tableaux.
  \item The \emph{insertion tableau}~$P$ is a semistandard
    Young tableau with entries in $\{1, \ldots, n\}$.
    In particular, $\max(P) = \max(\pi) = n$.
  \item The \emph{recording tableau}~$Q$ encodes the order
    in which cells are filled during the growth process.
  \item The shape of~$P$ (= shape of~$Q$) is determined by
    the interlacing structure.
\end{itemize}

\begin{remark}
For cylindric partitions, the cyclic boundary condition means
the growth diagram ``wraps around.'' This is the setting of
the cylindric growth diagrams developed
in~\cite{langer-thesis,langer-part-II}.
The RSK output is no longer a pair of tableaux but rather
a pair of \emph{cylindric tableaux}, and the analysis
is more delicate.
\end{remark}

\subsection{What max entry controls in RSK}

The maximum entry~$n$ governs three equivalent aspects of
the RSK decomposition:

\begin{center}
\begin{tabular}{l|l}
  \textbf{Model} & \textbf{Max entry $= n$ means} \\
  \hline
  Level sets & $n$ nested layers \\
  NILP & $n$ non-intersecting paths \\
  Growth diagram & $n$ rows (insertion steps) \\
  RSK output & $\max(P) = n$ \\
  Schur functions & alphabet of size~$n$
\end{tabular}
\end{center}

\begin{remark}[What max entry does \emph{not} control]
Under RSK, the maximum entry of the input does \emph{not}
determine:
\begin{itemize}
  \item The shape of the output tableaux (which depends on
    the full interlacing structure, not just the depth);
  \item The number of rows or columns of the output
    (Schensted's theorem relates the longest increasing
    subsequence to the first row, not to the maximum entry);
  \item Any clean statistic on the \emph{recording}
    tableau~$Q$.
\end{itemize}
Under the Hillman--Grassl correspondence~\cite{hillman-grassl}
(a different bijection for RPPs), the max entry of the RPP
has no known clean characterization on the output side.
\end{remark}

%% ========================================================================
\section{Gansner's trace generating function}
\label{sec:gansner}
%% ========================================================================

Gansner~\cite{gansner-1981} studied a different multivariate
refinement of RPPs: the \emph{trace generating function},
where one tracks the sum of entries along each anti-diagonal.

For an RPP~$\pi$ of shape~$\lambda$, define the $d$-th trace as
$t_d(\pi) = \sum_{j-i=d} \pi_{ij}$. Gansner showed:
\begin{equation}\label{eq:gansner}
  \sum_{\pi \in \RPP(\lambda)}
  \prod_{d} x_d^{t_d(\pi)}
  = \prod_{u = (i,j) \in \lambda}
  \frac{1}{1 - x_{j-i}^{h_\lambda(u)}}.
\end{equation}
Under the specialization $x_d \mapsto q$ for all $d$,
this recovers Stanley's formula~\eqref{eq:stanley-hook}.

\begin{remark}[Gansner vs.\ frequency refinement]
Gansner's trace GF and our frequency GF
(Section~\ref{sec:frequency}) are \emph{different}
multivariate refinements:
\begin{center}
\begin{tabular}{l|l|l}
  & \textbf{Gansner (traces)} &
  \textbf{Frequency} \\
  \hline
  Variable & $x_d$ per diagonal $d$ &
  $z_k$ per value $k$ \\
  Statistic & $t_d = \sum_{j-i=d} \pi_{ij}$ &
  $\#_k = |\{(i,j) : \pi_{ij} = k\}|$ \\
  Symmetry & Symmetric in each $x_d$
  & Quasi-symmetric in $z_k$'s \\
  Product formula & Yes \eqref{eq:gansner} & No \\
  Max entry & Not directly visible &
  $\max\{k : \#_k > 0\}$
\end{tabular}
\end{center}
The frequency refinement can ``see'' the max entry
(it is $\max\{k : z_k \text{ appears}\}$), while Gansner's
trace refinement cannot (the max entry is not a linear
function of the diagonal sums).
\end{remark}

%% ========================================================================
\section{The E--H factorization and Warnaar's $Q$ polynomials}
\label{sec:EH}
%% ========================================================================

\subsection{Extracting $Q_{n,c}(q)$}

The \emph{Euler--Hamburger factorization} separates
$\GK_c(z,q)$ into a universal factor and a profile-dependent
remainder:
\begin{equation}\label{eq:EH}
  \GK_c(z,q) = \frac{1}{\qpoch{zq}{q}} \cdot F_c(z,q),
  \qquad
  F_c(z,q) = \qpoch{zq}{q} \cdot \GK_c(z,q).
\end{equation}

Writing $F_c(z,q) = \sum_{n \geq 0} h_n(q) z^n$, the
$Q$ polynomials are:
\begin{equation}\label{eq:Q-def}
  Q_{n,c}(q) = \qpochn{q}{q}{n} \cdot h_n(q)
  = \qpochn{q}{q}{n}
  \cdot [z^n]\bigl(\qpoch{zq}{q} \cdot \GK_c(z,q)\bigr).
\end{equation}

\begin{conjecture}[Warnaar~\cite{warnaar-2021}]
\label{conj:warnaar}
For all $n \geq 0$ and all profiles~$c$ with
$d = |c| \not\equiv 0 \pmod{3}$, the polynomial
$Q_{n,c}(q)$ has nonneg integer coefficients.
\end{conjecture}

\subsection{The max-entry structure of $Q_n$}

The extraction~\eqref{eq:Q-def} can be rewritten
using the ``ground set'' formulation:
\begin{equation}\label{eq:Q-ground-set}
  Q_{n,c}(q) = \sum_{j=0}^{n}\;
  \sum_{\substack{S \subseteq \{1,\ldots,n\} \\ |S|=j}}\;
  \sum_{U \subseteq \{1,\ldots,n-j\}}\;
  \sum_{\substack{\nu \in \CPP(c) \\
  \max(\nu) = n-j}}
  (-1)^{|S|+|U|}\, q^{\sigma(S)+\sigma(U)+|\nu|}.
\end{equation}

This makes explicit that $Q_n$ involves CPPs at
\emph{every} max entry from $0$ to $n$: the $j$-th term
uses CPPs with $\max = n - j$. The alternating sign
$(-1)^{|S|+|U|}$ must produce miraculous cancellation to
yield a nonneg polynomial.

\subsection{The $q$-binomial decomposition}

The relationship between $Q_n$ and the ``corrected''
coefficients $h_m = \qpochn{q}{q}{m} \cdot g_m$ is:
\begin{equation}\label{eq:q-binom-decomp}
  Q_n = \sum_{j=0}^{n} (-1)^j\, q^{\binom{j+1}{2}}\,
  \qbin{n}{j}\, h_{n-j},
\end{equation}
with inverse:
\begin{equation}\label{eq:q-binom-inverse}
  h_m = \sum_{j=0}^{m} q^j\, \qbin{m}{j}\, Q_{m-j}.
\end{equation}

The polynomials $h_m(q)$ are nonneg for all tested
profiles (a key finding), and the alternating
$q$-binomial sum~\eqref{eq:q-binom-decomp} must preserve
nonnegativity --- this is the core mystery.

%% ========================================================================
\section{Frequency refinement at fixed max entry}
\label{sec:frequency}
%% ========================================================================

\subsection{The frequency generating function}

For a CPP $\pi = (\nu^{(1)}, \ldots, \nu^{(k)})$, the
\emph{value frequency} of~$k$ is
$\#_k(\pi) = |\{(i,j) : \nu^{(i)}_j = k\}|$.
The \emph{frequency generating function} for all CPPs:
\[
  \widetilde{\GK}_c(\mathbf{z})
  = \sum_{\pi \in \CPP(c)}
  z_1^{\#_1(\pi)} z_2^{\#_2(\pi)} \cdots
\]
is quasi-symmetric and $L$-positive
(see~\cite{langer-qsym}).

At fixed max entry:
\[
  \tilde{g}_n = \sum_{\max(\pi) = n}
  z_1^{\#_1(\pi)} \cdots z_n^{\#_n(\pi)}
  \;\in\; \Z[z_1, \ldots, z_n].
\]

\subsection{Partial quasi-symmetry}

The function $\tilde{g}_n$ is \emph{not} quasi-symmetric
in $z_1, \ldots, z_n$ (the value~$n$ is privileged: every
CPP with $\max = n$ has $\#_n \geq 1$, but may have
$\#_1 = 0$). However, it retains a partial quasi-symmetry:

\begin{observation}[Computational]\label{obs:partial-qsym}
For all tested profiles ($d \leq 5$, $n \leq 3$): writing
\[
  \tilde{g}_n = \sum_{m \geq 1} z_n^m \cdot
  f_{n,m}(z_1, \ldots, z_{n-1}),
\]
each slice $f_{n,m}$ is quasi-symmetric in
$z_1, \ldots, z_{n-1}$. Equivalently,
$\tilde{g}_n \in z_n \cdot \QSym(z_1, \ldots, z_{n-1})[z_n]$.
\end{observation}

The argument: relabelling the non-$n$ values
$\{v_1 < \cdots < v_k\} \subset \{1, \ldots, n-1\}$ to
any other increasing sequence
$\{w_1 < \cdots < w_k\} \subset \{1, \ldots, n-1\}$
preserves the interlacing conditions, because:
\begin{enumerate}
  \item Interlacing among non-$n$ entries depends only
    on their relative order;
  \item Every non-$n$ entry $v < n$, so all inequalities
    with $n$-entries are automatically satisfied.
\end{enumerate}

\subsection{Failure of slice $L$-positivity}

While the full $\widetilde{\GK}_c$ is $L$-positive, the
individual slices $f_{n,m}$ are \emph{not}. For $n = 3$
(two non-$n$ variables), the $M$-coefficients satisfy
$a_{(k)} > a_{(j,k)}$ for all tested profiles, giving
$L_{(a,b)} = M_{(a,b)} - M_{(a+b)} < 0$.

\subsection{The variable-count mismatch}
\label{sec:mismatch}

The E--H factorization~\eqref{eq:EH} couples objects living
in different numbers of variables:
\begin{itemize}
  \item The factor $\qpochn{q}{q}{n} = \prod_{k=1}^{n}(1 - q^k)$
    lifts to the elementary symmetric functions
    $e_1, e_2, \ldots$ in \emph{infinitely many} variables
    (or equivalently, to $\prod_{k \geq 1}(1 - z_k)$
    in $\QSym$).
  \item The coefficient $\tilde{g}_n$ is a polynomial in
    \emph{only $n$ variables} $z_1, \ldots, z_n$.
\end{itemize}

The product $\qpochn{q}{q}{n} \cdot g_n(q)$ in one variable
makes sense (multiply two polynomials in~$q$), but any
multivariate lift must reconcile infinite-variable and
finite-variable objects. This is a structural obstacle to
extending the $\QSym$ framework from $\widetilde{\GK}_c$
(where it succeeds) to $Q_{n,c}$ (where it fails).

%% ========================================================================
\section{Summary: what is and isn't known}
\label{sec:summary}
%% ========================================================================

\begin{center}
\begin{tabular}{l|c|c}
  \textbf{Object / Property}
  & \textbf{RPPs}
  & \textbf{CPPs} \\
  \hline
  Max entry = number of layers
  & Classical & Classical \\
  Hook-product formula ($z = 1$)
  & Stanley~\cite{stanley-1972}
  & Borodin~\cite{borodin-2007} \\
  Max $\leq n$ via Schur functions
  & Eq.~\eqref{eq:rpp-schur}
  & Transfer matrix~\eqref{eq:transfer-matrix} \\
  Growth diagram / RSK
  & Fomin~\cite{fomin-1995}
  & Cylindric GD~\cite{langer-part-II} \\
  Trace GF (Gansner)
  & \cite{gansner-1981}
  & Not studied \\
  Frequency GF quasi-symmetric
  & Yes (P-partition)
  & Yes~\cite{langer-qsym} \\
  Frequency GF $L$-positive
  & Expected
  & Yes~\cite{langer-qsym} \\
  Partial QSym at fixed max
  & Expected
  & Yes (Obs.~\ref{obs:partial-qsym}) \\
  Slice $L$-positivity
  & Unknown
  & \textbf{No} (fails at $n \geq 3$) \\
  E--H factorization
  & ---
  & Warnaar~\cite{warnaar-2021} \\
  $Q_{n,c}(q) \geq 0$
  & ---
  & \textbf{Open} (Conj.~\ref{conj:warnaar})
\end{tabular}
\end{center}

\medskip

\noindent\textbf{Gaps in the literature:}
\begin{enumerate}
  \item The frequency refinement $\tilde{g}_n$ (at fixed
    max entry) does not appear to have been studied, for
    either RPPs or CPPs. The partial quasi-symmetry
    (Observation~\ref{obs:partial-qsym}) appears new.
  \item Gansner's trace generating function has not been
    extended to cylindric partitions. The cyclic boundary
    condition complicates the diagonal structure.
  \item The Hillman--Grassl image of the maximum entry is
    unknown. A clean characterization could yield a new
    approach to the hook-product formula.
  \item The failure of slice $L$-positivity (despite
    aggregate $L$-positivity) has not been observed before
    and may constrain which algebraic approaches can prove
    positivity of~$Q_{n,c}$.
\end{enumerate}

%% ========================================================================
%% References
%% ========================================================================

\begin{thebibliography}{99}

\bibitem{borodin-2007}
A.~Borodin,
Periodic Schur process and cylindric partitions,
\emph{Duke Math.\ J.}\ \textbf{140} (2007), 391--468.

\bibitem{fomin-1986}
S.~V.~Fomin,
Generalized Robinson--Schensted--Knuth correspondence,
\emph{Zap.\ Nauchn.\ Sem.\ LOMI}\ \textbf{155} (1986), 156--175.

\bibitem{fomin-1995}
S.~Fomin,
Schur operators and Knuth correspondences,
\emph{J.~Combin.\ Theory Ser.~A}\ \textbf{72} (1995), 277--292.

\bibitem{gansner-1981}
E.~R.~Gansner,
The Hillman--Grassl correspondence and the enumeration of
reverse plane partitions,
\emph{J.~Combin.\ Theory Ser.~A}\ \textbf{30} (1981), 71--89.

\bibitem{gessel-krattenthaler-1997}
I.~M.~Gessel and C.~Krattenthaler,
Cylindric partitions,
\emph{Trans.\ Amer.\ Math.\ Soc.}\ \textbf{349} (1997), 429--479.

\bibitem{hillman-grassl}
A.~P.~Hillman and R.~M.~Grassl,
Reverse plane partitions and tableau hook numbers,
\emph{J.~Combin.\ Theory Ser.~A}\ \textbf{21} (1976), 216--221.

\bibitem{langer-thesis}
R.~Langer,
\emph{Enumeration of cylindric plane partitions},
PhD thesis, Universit\'e Claude Bernard Lyon~1, 2021.

\bibitem{langer-part-II}
R.~Langer,
Enumeration of cylindric plane partitions~-- Part~II,
arXiv:1209.1807.

\bibitem{langer-qsym}
R.~Langer,
Quasi-symmetric refinement of the cylindric partition
generating function, 2025. In preparation.

\bibitem{stanley-1972}
R.~P.~Stanley,
Ordered structures and partitions,
\emph{Mem.\ Amer.\ Math.\ Soc.}\ \textbf{119}, 1972.

\bibitem{stanley-ec2}
R.~P.~Stanley,
\emph{Enumerative Combinatorics}, Volume~2,
Cambridge University Press, 1999.

\bibitem{warnaar-2021}
S.~O.~Warnaar,
The $A_2$ Andrews--Gordon identities and cylindric partitions,
arXiv:2111.07550.

\end{thebibliography}

\end{document}
